\documentclass[11pt]{article}
\usepackage[margin=1in]{geometry}
\usepackage{amsmath,amssymb}
\usepackage{bm}
\usepackage{booktabs}
\usepackage{xcolor}
\usepackage{hyperref}
\hypersetup{colorlinks=true,linkcolor=blue,urlcolor=blue}

\newcommand{\rhat}{\hat{\mathbf r}}
\newcommand{\thetahat}{\hat{\boldsymbol\theta}}
\newcommand{\phihat}{\hat{\boldsymbol\phi}}

\title{Electric field computation in \texttt{sourceField1d}}
\author{Anna Kelbert with Claude}
\date{September 2026}

\begin{document}
\maketitle

\begin{abstract}
\texttt{sourceField1d} (\texttt{EARTH/FWD/field1d.f90}) computes both the magnetic field
$\mathbf{H}$ (poloidal) and the electric field $\mathbf{E}$ (toroidal) of the 1D layered-sphere
solution, from a single scalar toroidal potential $T_l(r)$. $\mathbf{H}$ and $\mathbf{E}$ are
built from opposite vector spherical harmonic (VSH) patterns of the same potential, related by
Faraday's law. This note derives that relation, gives the exact formulas and grid staggering
\texttt{sourceField1d} implements, and records the one physical convention choice (radial
direction) that distinguishes this solver from its companion, \texttt{field1d\_s2.f90}.
\end{abstract}

\section{Poloidal and toroidal decomposition}

Any smooth vector field $\mathbf F$ on a spherical surface splits into three orthogonal parts:
\begin{equation}
\mathbf F = F_r\,\rhat \;+\; \nabla_a P \;+\; \rhat\times\nabla_a T .
\end{equation}
$\nabla_a P$ (the tangential gradient of a scalar $P$) is the \textbf{poloidal} part;
$\rhat\times\nabla_a T$ (the curl of a radial vector) is the \textbf{toroidal} part. For a
magnetotelluric source, $\mathbf H$ is poloidal ($H_r\ne0$) and $\mathbf E$ is toroidal ($E_r=0$,
purely tangential) --- the standard TE-mode picture of deep-Earth induction.

In vector spherical harmonics (degree $l$, order $m$), with \texttt{vsharm} computing
$Y(l,m{+}1)=Y_l^m$, $Y_t(l,m{+}1)=\partial Y/\partial\theta$,
$Y_p(l,m{+}1)=(\partial Y/\partial\phi)/\sin\theta$:
\begin{align}
\Psi_{lm} &= \nabla_a Y_{lm} \;\to\; (Y_t\,\thetahat + Y_p\,\phihat) &&\text{POLOIDAL pattern} \\
\Phi_{lm} &= \rhat\times\nabla_a Y_{lm} \;\to\; (-Y_p\,\thetahat + Y_t\,\phihat) &&\text{TOROIDAL pattern}
\end{align}
$\mathbf H$ uses $\Psi_{lm}$ (theta component $Y_t$, phi component $+Y_p$); $\mathbf E$ uses
$\Phi_{lm}$ (theta component $-Y_p$, phi component $+Y_t$) --- the same pair $(Y_t,Y_p)$, swapped
and one sign flipped. This is the fundamental Maxwell duality for the TE mode:
\begin{equation}
\nabla\times\mathbf E = i\omega\mu_0\mathbf H \ \text{(Faraday)}, \qquad
\nabla\times\mathbf H = \sigma\mathbf E \ \text{(Amp\`ere, quasi-static)} .
\end{equation}
Faraday's toroidal curl gives a poloidal field, and vice versa.

\section{The potentials, and the $H$ field}
\label{sec:H}

\texttt{sourcePotential} returns, at every radius and degree:

\begin{table}[h]
\centering
\begin{tabular}{llll}
\toprule
Array & Value & Radii & Drives \\
\midrule
\texttt{Tnr(k,l)}  & $T(r)$   & $R_r(1{:}N_r)$, cell centers & $H_r$, both $E$ components \\
\texttt{Tnsp(k,l)} & $dT/dr$  & $R_s(1{:}N_r{+}1)$, cell faces & $H_\theta$, $H_\phi$ \\
\bottomrule
\end{tabular}
\end{table}

\texttt{Tnsp = d(Tnr)/dr} --- the physical radial derivative of \texttt{Tnr}.

\texttt{sourceField1d} assembles, per component (Fortran coefficient ordering
$m=0,{+}1,{-}1,{+}2,{-}2,\dots$):
\begin{align}
H_x(i,j,k)\;[\phi]     &= \sum_l Y_p(l,\cdot)\,\mathtt{coefl}\,\mathtt{Tnsp}(k,l)\,R_0^2/R_s(k)/l(l{+}1) \\
H_y(i,j,k)\;[\theta]   &= \sum_l Y_t(l,\cdot)\,\mathtt{coefl}\,\mathtt{Tnsp}(k,l)\,R_0^2/R_s(k)/l(l{+}1) \\
H_z(i,j,k)\;[r]        &= \sum_l Y(l,\cdot)\,\mathtt{coefl}\,\mathtt{Tnr}(k,l)\,R_0^2/R_r(k)^2
\end{align}
($H_r$ carries no $l(l{+}1)$ denominator --- consistent with how \texttt{Tnr}/\texttt{Tnsp} are
normalized in \texttt{sourcePotential}.) Each $l$-sum is built directly from the $m=0$ term plus
each $+m/-m$ pair (\S\ref{sec:pairing}) --- no post-hoc conjugation of any kind.

\section{Electric field: derivation and formulas}
\label{sec:E}

Start from $\nabla\times\mathbf E = i\omega\mu_0\mathbf H$, $E_r=0$. The theta-component of
Faraday's law gives
\begin{equation}
-\frac{1}{r}\frac{d(rE_\phi)}{dr} = i\omega\mu_0 H_\theta
= i\omega\mu_0\,\mathtt{Tnsp}(r)\,\frac{R_0^2}{r\,l(l{+}1)}\,Y_t .
\end{equation}
Integrating ($\mathtt{Tnsp}=d\,\mathtt{Tnr}/dr$):
\begin{equation}
E_\phi = -i\omega\mu_0\,\mathtt{Tnr}(r)\,\frac{R_0^2}{r\,l(l{+}1)}\,Y_t .
\end{equation}
The phi-component gives, by the same steps,
\begin{equation}
E_\theta = +i\omega\mu_0\,\mathtt{Tnr}(r)\,\frac{R_0^2}{r\,l(l{+}1)}\,Y_p .
\end{equation}
$E_r=0$ throughout (TE mode). The radial component of Faraday's law is automatically satisfied
--- not an independent constraint --- since $E_\theta$/$E_\phi$ were built from the SAME potential
$T_l(r)$ whose value (\texttt{Tnr}) and derivative (\texttt{Tnsp}) already solve one consistent
radial ODE; the poloidal/toroidal formalism guarantees the two are compatible by construction.

\paragraph{Result:}
\begin{align}
E_\theta &= +i\omega\mu_0\,\mathtt{Tnr}(r)\,R_0^2/(r\,l(l{+}1))\;Y_p \\
E_\phi   &= -i\omega\mu_0\,\mathtt{Tnr}(r)\,R_0^2/(r\,l(l{+}1))\;Y_t \\
E_r      &= 0
\end{align}
$E$ uses the existing \texttt{Tnr} array directly --- no separate potential is needed.

\subsection*{Angular patterns compared}

\begin{table}[h]
\centering
\begin{tabular}{lllll}
\toprule
Component & $H$ pattern & $H$ potential & $E$ pattern & $E$ potential \\
\midrule
theta ($H_y$/$E_y$) & $+Y_t$ & \texttt{Tnsp} at $R_s$ & $+Y_p$ & $i\omega\mu_0\,\mathtt{Tnr}$ at $R_r$ \\
phi ($H_x$/$E_x$)   & $+Y_p$ & \texttt{Tnsp} at $R_s$ & $-Y_t$ & $i\omega\mu_0\,\mathtt{Tnr}$ at $R_r$ \\
radial ($H_z$/$E_z$) & $+Y$  & \texttt{Tnr} at $R_r$   & $0$    & --- \\
\bottomrule
\end{tabular}
\end{table}

$E_{\rm tang}$ pairs with $H_r$ (same potential $T(r)$, different angular pattern); $H_{\rm tang}$
pairs with $dT/dr$ (\texttt{Tnsp}). This is the Faraday-Amp\`ere duality.

\section{Coefficient assembly: the $\pm m$ pairing}
\label{sec:pairing}

Fortran coefficient ordering differs from MATLAB's:

\begin{table}[h]
\centering
\begin{tabular}{lll}
\toprule
Physical $m$ & Fortran \texttt{coefl} index & MATLAB \texttt{amp} index \\
\midrule
$m=0$   & $1$     & $1$ \\
$m=+m$  & $2m$    & $m{+}1$ \\
$m=-m$  & $2m{+}1$ & $l{+}1{+}m$ \\
\bottomrule
\end{tabular}
\end{table}

Each degree-$l$ sum adds the $m=0$ term directly, then each $+m/-m$ pair via the Condon-Shortley
identity $Y_l^{-m}=(-1)^m\overline{Y_l^m}$:
\begin{equation}
C = \frac{Y(l,m{+}1)\,\mathtt{coefl}(2m) + (-1)^m\overline{Y(l,m{+}1)}\,\mathtt{coefl}(2m{+}1)}{l(l{+}1)}
\end{equation}
for the $l(l{+}1)$-normalized components ($H_\theta$, $H_\phi$, $E_\theta$, $E_\phi$); $H_r$ uses
the same pairing without the $l(l{+}1)$ divisor. $E_x$'s accumulated pair carries an overall
leading minus sign, matching the $-Y_t$ pattern above. The $(-1)^m$ factor is required: omitting
it reproduces the physical sum only for even $m$.

No component is conjugated after assembly --- each $l$-sum, once the $\pm m$ pairing above is
applied term-by-term, is already the correct physical field.

\section{Radial (source-amplitude) normalization}
\label{sec:radial}

The formulas in \S\ref{sec:H}--\ref{sec:E} are written directly in terms of \texttt{coefl} ---
the coefficient array as consumed natively by \texttt{sourceField1d} --- and
\texttt{Tnr}/\texttt{Tnsp}, the radial potential/derivative solved once per degree $l$,
independent of $m$ and of the coefficient values. Because \texttt{Tnr}/\texttt{Tnsp} never depend
on \texttt{coeff}, every formula above ($H_r$, $H_\theta$, $H_\phi$, $E_\theta$, $E_\phi$) is
exactly \textbf{linear in \texttt{coefl}} --- a coefficient's absolute scale enters every field
component identically.

\texttt{field1d.f90} natively expects $\mathtt{coeff}(l,m)=1$ to mean a \textbf{unit external
surface radial field}: $H_r(R_0) = Y_l^m(\theta,\phi)$ --- the \texttt{RADIAL\_SURFACE} convention
in \texttt{EARTH/FWD/output\_convention.f90}. Three other conventions are supported for source
coefficients authored differently: an amplitude of the external potential term
$T_l^{\rm ext}(r)=\alpha_l^T(r/R_0)^{l+1}$ (\texttt{RADIAL\_MULTIPOLE}, \texttt{field1d\_s2.f90}'s
own native convention, unit $\alpha_l^T=1$); a literal $r^{l+1}$ coefficient
(\texttt{RADIAL\_DIMENSIONAL}, Sun \& Egbert's own text convention, $r$ in metres); or the
classical external geomagnetic potential's Gauss coefficient (\texttt{RADIAL\_POTENTIAL},
$V=R_0\sum_l(r/R_0)^l\varepsilon_l^m Y_l^m$, used for ionospheric-current/Sq sources).

To use coefficients authored in one of these other conventions with the formulas in
\S\ref{sec:H}--\ref{sec:E}, rescale them \textbf{per degree $l$} (the same factor for every $m$
of that degree, before substituting as \texttt{coefl}):
\begin{equation}
\mathtt{coefl\_for\_formula}(l,m) = \mathtt{coeff\_in\_other\_convention}(l,m)\cdot
\frac{A(\mathtt{SURFACE},l)}{A(\mathtt{other\_convention},l)}
\end{equation}
with the per-degree amplitude $A(l)$, relative to the \texttt{MULTIPOLE} baseline (identical to
\texttt{output\_convention.f90}'s \texttt{radial\_amplitude()}):

\begin{table}[h]
\centering
\begin{tabular}{ll}
\toprule
\texttt{radial\_norm} & $A(l)$ \\
\midrule
\texttt{MULTIPOLE}   & $1$ \\
\texttt{SURFACE}     & $R_0^2/(l(l{+}1))$ \\
\texttt{DIMENSIONAL} & $R_0^{l+1}$ (overflows double precision for $l\gtrsim44$) \\
\texttt{POTENTIAL}   & $R_0^2/(l{+}1)$ \\
\bottomrule
\end{tabular}
\end{table}

Because this rescale is a single real scalar per degree $l$, applied identically wherever
\texttt{coefl} appears in \S\ref{sec:H}--\ref{sec:E}, it multiplies $H_r$, $H_\theta$, $H_\phi$,
$E_\theta$, $E_\phi$ all by the \textbf{same} factor for that degree --- an overall amplitude
change only. It never touches the angular pattern ($Y$/$Y_t$/$Y_p$), the relative $H$/$E$
component ratios, or the impedance $Z=E/H$ (amplitude-independent by construction). This
linearity is exactly why rescaling the source coefficients before the solve
(\texttt{rescale\_source\_radial}) and rescaling the finished $H$/$E$ output by the same factor
afterward are equivalent operations.

\section{Grid staggering}

$H_r$/$E$'s \texttt{Tnr} live at cell-center radii $R_r$; $H_\theta$/$H_\phi$'s \texttt{Tnsp} at
cell-face radii $R_s$ ($k$ loops over $N_{rr}=N_r$ vs.\ $N_{rs}=N_r{+}1$ respectively).

\begin{table}[h]
\centering
\begin{tabular}{lllll}
\toprule
Component & theta eval & phi eval & radii & $j$ range \\
\midrule
$H_x$ (phi)    & node $\theta(j)$      & mid $\phi{+}d\phi/2$ & $R_s$ & $2..N_t$ \\
$H_y$ (theta)  & mid $\theta{+}d\theta/2$ & node $\phi(i)$    & $R_s$ & $1..N_t$ \\
$H_z$ (r)      & node $\theta(j)$      & node $\phi(i)$       & $R_r$ & $j_1..j_2$ (pole-aware) \\
$E_y$ (theta)  & node $\theta(j)$      & mid $\phi{+}d\phi/2$ & $R_r$ & $j_1..j_2$ (pole-aware) \\
$E_x$ (phi)    & mid $\theta{+}d\theta/2$ & node $\phi(i)$    & $R_r$ & $1..N_t$ \\
$E_z$ (r)      & --- & --- & --- & zero \\
\bottomrule
\end{tabular}
\end{table}

Each $E$ component shares its angular stagger with the $H$ component using the \textbf{same VSH
factor} --- $E_y$ ($Y_p$) with $H_x$ (also $Y_p$, node theta/mid phi), $E_x$ ($Y_t$) with $H_y$
(also $Y_t$, mid theta/node phi) --- not with the $H$ component it pairs with physically via
Faraday's law. $E_y$, not $E_x$, is therefore the pole-adjacent (node-theta) component and
carries the pole-aware $j_1,j_2$ range: $j_1=1$ unless $\theta(1)$ is a true pole ($j_1=2$ then);
$j_2=N_t{+}1$ unless $\theta(N_t{+}1)$ is a true pole ($j_2=N_t$ then). This range is only ever
narrowed for a genuinely global grid reaching both poles --- a regional grid keeps its full node
range. \texttt{node\_P\_lm} (shared with $H_z$) is reused for $E_y$; \texttt{edge\_P\_lm}
(shared with $H_y$) is reused for $E_x$ --- no separate Legendre computation is needed for $E$.

\section{Scope and the radial-direction convention}

$E_r=0$ is exact for the TE (poloidal-$H$) source this solver targets; a source with TM-mode
content would need separate treatment (not needed for standard Sq/ring-current MT sources).

\texttt{field1d.f90}'s raw output above is natively $e^{-i\omega t}$, with $\rhat$ pointing away
from Earth's center (increasing radius) --- identical, in every bookkeeping respect but one, to
the companion solver \texttt{field1d\_s2.f90}. This was a deliberate choice, not the solver's
original one: \texttt{field1d.f90} was written in 2011 as a faithful implementation of Kelbert
(2006)'s own conventions (the Global3D Fortran code in ModEM), which use $\rhat$ pointing toward
Earth's center and $e^{+i\omega t}$ output --- genuinely different, equally valid physical
conventions. \texttt{field1d.f90} was later re-pointed to match \texttt{field1d\_s2.f90} so the
two independently-derived solvers agree natively with no compensating sign in comparison code;
Kelbert (2006)'s original convention remains available as the named \texttt{KELBERT2006} preset
in \texttt{EARTH/FWD/output\_convention.f90}. See \texttt{docs/output\_conventions.pdf} for the
full convention system and the derivation of exactly which field components that
radial-direction choice affects.

\end{document}
